\chapter{2003年上海卷（春）}
\section{填空题}
\begin{problem}
已知函数 \( f(x) = \sqrt{x} + 1 \)，则 \( f^{-1}(3) = \)\fillinblank{}.
\end{problem}

\begin{answer}
\(4\)
\end{answer}

\begin{solution}
令 \(f(x)=3\)，则 \(\sqrt{x}+1=3\)，所以 \(\sqrt{x}=2\)，从而 \(x=4\). 因此 \(f^{-1}(3)=4\).
\end{solution}

\begin{problem}
直线 \( y = 1 \) 与直线 \( y = \sqrt{3}x + 3 \) 的夹角为\fillinblank{}.
\end{problem}

\begin{answer}
\(\dfrac{\pi}{3}\)
\end{answer}

\begin{solution}
直线 \(y=1\) 的斜率为 \(0\)，直线 \(y=\sqrt{3}x+3\) 的斜率为 \(\sqrt{3}\). 设两直线的锐角夹角为 \(\theta\)，则
\[
\tan\theta=\left|\frac{\sqrt{3}-0}{1+0\cdot\sqrt{3}}\right|=\sqrt{3}.
\]
因为 \(0<\theta<\dfrac{\pi}{2}\)，所以 \(\theta=\dfrac{\pi}{3}\).
\end{solution}

\begin{problem}
已知点 \(P(\tan \alpha,\cos \alpha)\) 在第三象限，则角 \(\alpha\) 的终边在第\fillinblank{}象限.
\end{problem}

\begin{answer}
二
\end{answer}

\begin{solution}
点 \(P\) 在第三象限，所以 \(\tan\alpha<0\)，且 \(\cos\alpha<0\). \(\cos\alpha<0\) 表明角 \(\alpha\) 的终边在第二或第三象限，\(\tan\alpha<0\) 表明其终边在第二或第四象限. 两者的交集是第二象限，因此角 \(\alpha\) 的终边在第二象限.
\end{solution}

\begin{problem}
直线 \( y = x - 1 \) 被抛物线 \( y^{2} = 4x \) 截得线段的中点坐标是\fillinblank{}.
\end{problem}

\begin{answer}
\((3,2)\)
\end{answer}

\begin{solution}
由 \(y=x-1\)，得 \(x=y+1\). 代入抛物线方程，得 \(y^{2}=4(y+1)\)，即 \(y^{2}-4y-4=0\). 因此 \(y=2\pm2\sqrt{2}\)，相应地 \(x=3\pm2\sqrt{2}\).

两交点的中点为
\[
\left(\frac{(3+2\sqrt{2})+(3-2\sqrt{2})}{2},\frac{(2+2\sqrt{2})+(2-2\sqrt{2})}{2}\right)=(3,2).
\]
\end{solution}

\begin{problem}
已知集合 \(A = \{x \mid |x| \leqslant 2,\ x \in \R\}\)，\(B = \{x \mid x \geqslant a\}\)，且 \(A \subseteq B\)，则实数 \(a\) 的取值范围是\fillinblank{}.
\end{problem}

\begin{answer}
\((-\infty,-2]\)
\end{answer}

\begin{solution}
由 \(|x|\leqslant2\) 可知 \(A=[-2,2]\)，而 \(B=[a,+\infty)\). 要使 \(A\subseteq B\)，只需使 \(A\) 的最小值 \(-2\) 不小于 \(B\) 的下界 \(a\)，即 \(a\leqslant-2\). 所以 \(a\) 的取值范围是 \((-\infty,-2]\).
\end{solution}

\begin{problem}
已知 \(z\) 为复数，则 \(z + \overline{z} > 2\) 的一个充要条件是 \(z\) 满足\fillinblank{}.
\end{problem}

\begin{answer}
\(z=u+\i v\)，其中 \(u>1\)，\(u,v\in\R\)
\end{answer}

\begin{solution}
设 \(z=u+\i v\)，其中 \(u,v\in\R\)，则 \(\overline{z}=u-\i v\)，从而 \(z+\overline{z}=2u\). 因此
\[
z+\overline{z}>2\Longleftrightarrow2u>2\Longleftrightarrow u>1.
\]
所以一个充要条件是 \(z=u+\i v\)，其中 \(u,v\in\R\) 且 \(u>1\).
\end{solution}

\begin{problem}
若过两点 \(A(-1,0)\)，\(B(0,2)\) 的直线 \(l\) 与圆 \((x-1)^{2}+(y-a)^{2}=1\) 相切，则 \(a=\)\fillinblank{}.
\end{problem}

\begin{answer}
\(4-\sqrt{5}\) 或 \(4+\sqrt{5}\)
\end{answer}

\begin{solution}
直线 \(l\) 的斜率为 \(2\)，所以其方程为 \(y=2x+2\)，即 \(2x-y+2=0\). 圆心为 \(C(1,a)\)，半径为 \(1\).

由直线与圆相切，圆心到直线的距离等于半径，故
\[
\frac{|2-a+2|}{\sqrt{2^{2}+(-1)^{2}}}=1,
\]
即 \(|4-a|=\sqrt{5}\). 所以 \(a=4-\sqrt{5}\) 或 \(a=4+\sqrt{5}\).
\end{solution}

\begin{problem}
不等式 \((\lg 20)^{2\cos x} > 1\)（\(x \in (0,\pi)\)）的解为\fillinblank{}.
\end{problem}

\begin{answer}
\(\left(0,\dfrac{\pi}{2}\right)\)
\end{answer}

\begin{solution}
因为 \(\lg20>1\)，所以对于实指数，有
\[
(\lg20)^{2\cos x}>1\Longleftrightarrow2\cos x>0\Longleftrightarrow\cos x>0.
\]
在 \(x\in(0,\pi)\) 内，\(\cos x>0\) 的解为 \(x\in\left(0,\dfrac{\pi}{2}\right)\).
\end{solution}

\begin{problem}
\(8\) 名世界网球顶级选手在上海大师赛上分成两组，每组各 \(4\) 人，分别进行单循环赛，每组决出前两名，再由每组的第一名与另一组的第二名进行淘汰赛，获胜者角逐冠、亚军，败者角逐第三、四名，则该大师赛共有\fillinblank{}\(\text{场}\)比赛.
\end{problem}

\begin{answer}
\(16\)
\end{answer}

\begin{solution}
每组 \(4\) 人进行单循环赛，每两人比赛一次，所以每组有 \(\mathrm C_{4}^{2}=6\) 场，两组共有 \(12\) 场.

两组的第一名分别与另一组的第二名进行淘汰赛，共有 \(2\) 场；两场比赛的胜者争夺冠、亚军，败者争夺第三、四名，又有 \(2\) 场. 因此比赛总场数为 \(12+2+2=16\) 场.
\end{solution}

\begin{problem}
若正三棱锥底面边长为 \(4\)，体积为 \(1\)，则侧面和底面所成二面角的大小等于\fillinblank{}（结果用反三角函数值表示）.
\end{problem}

\begin{answer}
\(\arctan\dfrac{3}{8}\)
\end{answer}

\begin{solution}
设正三棱锥的顶点为 \(S\)，底面中心为 \(O\)，底面一边的中点为 \(M\)，高为 \(SO=h\). 底面正三角形的面积为 \(4\sqrt{3}\)，由体积公式得
\[
1=\frac{1}{3}\cdot4\sqrt{3}\cdot h,
\]
所以 \(h=\dfrac{3}{4\sqrt{3}}=\dfrac{\sqrt{3}}{4}\).

在正三角形中，中心到边的距离为其内切圆半径，故 \(OM=\dfrac{2\sqrt{3}}{3}\). 过 \(SO\) 和 \(SM\) 作截面，截面垂直于棱 \(BC\)，其中 \(SM\) 与 \(OM\) 的夹角就是侧面与底面所成的二面角 \(\theta\). 在直角三角形 \(SOM\) 中，
\[
\tan\theta=\frac{SO}{OM}=\frac{\sqrt{3}/4}{2\sqrt{3}/3}=\frac{3}{8}.
\]
所以 \(\theta=\arctan\dfrac{3}{8}\).
\end{solution}

\begin{problem}
若函数 \( y = x^{2} + (a + 2)x + 3 \)，\( x \in [a,b] \) 的图象关于直线 \( x = 1 \) 对称，则 \( b = \)\fillinblank{}.
\end{problem}

\begin{answer}
\(6\)
\end{answer}

\begin{solution}
抛物线的对称轴为
\[
x=-\frac{a+2}{2}.
\]
题设图象关于 \(x=1\) 对称，所以 \(-\dfrac{a+2}{2}=1\)，得 \(a=-4\).

函数的定义区间为 \([a,b]\)，其图象关于 \(x=1\) 对称时，区间的两个端点也关于 \(x=1\) 对称，因此 \(\dfrac{a+b}{2}=1\). 代入 \(a=-4\)，得 \(b=6\).
\end{solution}

\begin{problem}
设 \( f(x) = \frac{1}{2^{x} + \sqrt{2}} \)，利用课本中推导等差数列前 \(n\) 项和公式的方法，可求得 \( f(-5) + f(-4) + \cdots + f(0) + \cdots + f(5) + f(6) \) 的值为\fillinblank{}.
\end{problem}

\begin{answer}
\(3\sqrt{2}\)
\end{answer}

\begin{solution}
对任意实数 \(x\)，令 \(t=2^{x-1/2}>0\)，则
\[
f(x)=\frac{1}{\sqrt{2}(t+1)},\qquad f(1-x)=\frac{1}{\sqrt{2}(1+t^{-1})}=\frac{t}{\sqrt{2}(t+1)}.
\]
所以 \(f(x)+f(1-x)=\dfrac{1}{\sqrt{2}}\).

将题中各项配对为 \(f(-5)+f(6)\)，\(f(-4)+f(5)\)，\(f(-3)+f(4)\)，\(f(-2)+f(3)\)，\(f(-1)+f(2)\)，\(f(0)+f(1)\)，共有 \(6\) 对，每对的和均为 \(\dfrac{1}{\sqrt{2}}\). 因此原式为 \(\dfrac{6}{\sqrt{2}}=3\sqrt{2}\).
\end{solution}

\section{单选题}
\begin{problem}
关于直线\(a\)，\(b\)，\(l\) 以及平面 \(M\)，\(N\)，下列命题中正确的是（\quad）

\choices
    {若 \(a \parallel M,\ b \parallel M\)，则 \(a \parallel b\)}
    {若 \(a \parallel M\)，\(b \perp a\)，则 \(b \perp M\)}
    {若 \(a \subset M\)，\(b \subset M\)，且 \(l \perp a\)，\(l \perp b\)，则 \(l \perp M\)}
    {若 \(a \perp M\)，\(a \parallel N\)，则 \(M \perp N\)}
\end{problem}

\begin{answer}
D
\end{answer}

\begin{solution}
选项 A 中，两条分别平行于同一平面的直线不一定平行；选项 B 中，直线 \(b\) 垂直于直线 \(a\) 不能推出 \(b\) 垂直于平面 \(M\)；选项 C 中，若 \(a\) 与 \(b\) 平行，直线 \(l\) 可能在平面 \(M\) 内同时垂直于它们，因而条件不足. 对于选项 D，直线 \(a\) 垂直于平面 \(M\)，所以 \(a\) 的方向为平面 \(M\) 的法向方向；又因 \(a\parallel N\)，该法向方向平行于平面 \(N\)，所以平面 \(N\) 与平面 \(M\) 垂直. 故选 D.
\end{solution}

\begin{problem}
复数 \(z = \frac{m - 2\i}{1 + 2\i}\)（\(m \in \R\)，\(\i\) 为虚数单位）在复平面上对应的点不可能位于（\quad）

\choices
    {第一象限}
    {第二象限}
    {第三象限}
    {第四象限}
\end{problem}

\begin{answer}
A
\end{answer}

\begin{solution}
分子分母同乘以 \(1-2\i\)，得
\[
z=\frac{(m-2\i)(1-2\i)}{5}=\frac{m-4}{5}-\frac{2m+2}{5}\i.
\]
因此对应点的横坐标为 \(\dfrac{m-4}{5}\)，纵坐标为 \(-\dfrac{2m+2}{5}\). 若点在第一象限，则需同时满足 \(m>4\) 和 \(m<-1\)，这是不可能的. 取 \(m<-1\)、\(-1<m<4\)、\(m>4\) 时，分别可得到第二、第三、第四象限的点. 故选 A.
\end{solution}

\begin{problem}
把曲线\(y \cos x + 2y - 1 = 0\) 先沿 \(x\) 轴向右平移 \(\frac{\pi}{2}\) 个单位，再沿 \(y\) 轴向下平移一个单位，得到的曲线方程是（\quad）

\choices
    {\((1 - y)\sin x + 2y - 3 = 0\)}
    {\((y - 1)\sin x + 2y - 3 = 0\)}
    {\((y + 1)\sin x + 2y + 1 = 0\)}
    {\(-(y + 1)\sin x + 2y + 1 = 0\)}
\end{problem}

\begin{answer}
C
\end{answer}

\begin{solution}
沿 \(x\) 轴向右平移 \(\dfrac{\pi}{2}\) 后，新坐标仍记为 \((x,y)\)，原方程中的 \(x\) 应替换为 \(x-\dfrac{\pi}{2}\)，所以得到 \(y\sin x+2y-1=0\). 再沿 \(y\) 轴向下平移一个单位，原方程中的 \(y\) 应替换为 \(y+1\)，于是
\[
(y+1)\sin x+2(y+1)-1=0,
\]
即 \((y+1)\sin x+2y+1=0\). 故选 C.
\end{solution}

\begin{problem}
关于函数 \( f(x) = (\sin x)^2 - \left(\frac{2}{3}\right)^{|x|} + \frac{1}{2} \)，有下面四个结论：

\begin{circlelist}
\item \( f(x) \) 是奇函数；
\item 当 \(x > 2003\) 时，\(f(x) > \frac{1}{2}\) 恒成立；
\item \( f(x) \) 的最大值是 \( \frac{3}{2} \)；
\item \( f(x) \) 的最小值是 \( -\frac{1}{2} \).
\end{circlelist}

其中正确结论的个数为（\quad）

\choices
    {\(1\) 个}
    {\(2\) 个}
    {\(3\) 个}
    {\(4\) 个}
\end{problem}

\begin{answer}
A
\end{answer}

\begin{solution}
因为 \((\sin x)^2\) 和 \(\left(\dfrac{2}{3}\right)^{|x|}\) 都是偶函数，所以 \(f(x)\) 是偶函数而不是奇函数，第一条结论错误. 当 \(x=k\pi>2003\) 时，\(\sin^2x=0\)，从而 \(f(x)=\dfrac{1}{2}-\left(\dfrac{2}{3}\right)^x<\dfrac{1}{2}\)，第二条结论错误.

又因为 \(0\leqslant\sin^2x\leqslant1\)，且 \(\left(\dfrac{2}{3}\right)^{|x|}>0\)，所以 \(f(x)<\dfrac{3}{2}\). 在 \(x=\dfrac{\pi}{2}+k\pi\) 处，\(f(x)=\dfrac{3}{2}-\left(\dfrac{2}{3}\right)^{|x|}<\dfrac{3}{2}\)，并且当 \(|x|\) 充分大时函数值可无限接近 \(\dfrac{3}{2}\)，故 \(\dfrac{3}{2}\) 不是函数的最大值，第三条结论错误.

最后，\(f(x)\geqslant0-1+\dfrac{1}{2}=-\dfrac{1}{2}\)，且 \(f(0)=-\dfrac{1}{2}\)，所以第四条结论正确. 正确结论只有 \(1\) 个，故选 A.
\end{solution}

\section{解答题}
\begin{problem}
解不等式组：\(\begin{cases}
x^{2} - 6x + 8 > 0,\\
\dfrac{x + 3}{x - 1} > 2.
\end{cases}\)
\end{problem}

\begin{solution}
由 \(x^{2}-6x+8>0\)，得 \((x-2)(x-4)>0\)，所以 \(x<2\) 或 \(x>4\).

又因为 \(x\neq1\)，且
\[
\frac{x+3}{x-1}-2=\frac{5-x}{x-1}>0,
\]
所以 \(1<x<5\). 两个不等式的解集取交集，得
\[
x\in(1,2)\cup(4,5).
\]
\end{solution}

\begin{problem}
已知函数 \( f(x) = A \sin (\omega x + \varphi) \)，（\(A > 0\)，\(\omega > 0\)，\(x \in \R\)）在一个周期内的图象如图所示，求直线 \( y = \sqrt{3} \) 与函数 \( f(x) \) 图象的所有交点的坐标.

\bitmapfigure[max width=0.66\linewidth]{img/2003/shanghai_spring/q18_fig1.png}
\end{problem}

\begin{solution}
由图象可知，函数的最大值为 \(2\)，最小值为 \(-2\)，所以 \(A=2\). 从最大值点 \(x=\dfrac{\pi}{2}\) 到最小值点 \(x=\dfrac{5\pi}{2}\) 的距离为半个周期，故周期为 \(4\pi\)，从而 \(\omega=\dfrac{1}{2}\). 图象在 \(x=-\dfrac{\pi}{2}\) 处过原点且向上，故可取
\[
f(x)=2\sin\left(\frac{x}{2}+\frac{\pi}{4}\right).
\]

令 \(f(x)=\sqrt{3}\)，得
\[
\sin\left(\frac{x}{2}+\frac{\pi}{4}\right)=\frac{\sqrt{3}}{2}.
\]
因此
\[
\frac{x}{2}+\frac{\pi}{4}=\frac{\pi}{3}+2k\pi\quad\text{或}\quad\frac{2\pi}{3}+2k\pi,
\qquad k\in\Z.
\]
解得 \(x=\dfrac{\pi}{6}+4k\pi\) 或 \(x=\dfrac{5\pi}{6}+4k\pi\). 所有交点的坐标为
\[
\left(\frac{\pi}{6}+4k\pi,\sqrt{3}\right)\quad\text{或}\quad\left(\frac{5\pi}{6}+4k\pi,\sqrt{3}\right),\qquad k\in\Z.
\]
\end{solution}

\begin{problem}
已知三棱柱 \(ABC - A_{1}B_{1}C_{1}\)，在某个空间直角坐标系中，\(\overrightarrow{AB}=\begin{pmatrix}
\frac{m}{2},\\
-\frac{\sqrt{3}m}{2},\\
0
\end{pmatrix}\)，\(\overrightarrow{AC}=\left\{m,0,0\right\}\)，\(\overrightarrow{AA_1}=\left\{0,0,n\right\}\)，其中 \(m,n>0\).

\begin{enumerate}
\item 证明：三棱柱 \(ABC - A_{1}B_{1}C_{1}\) 是正三棱柱；
\item 若 \(m = \sqrt{2}n\)，求直线 \(CA_{1}\) 与平面 \(A_{1}ABB_{1}\) 所成角的大小.
\end{enumerate}

\bitmapinlinefigure[0.33\linewidth][max width=0.33\linewidth]{img/2003/shanghai_spring/q19_fig1.png}
\end{problem}

\begin{solution}
\begin{enumerate}
\item 由已知向量可得 \(|AB|=m\)，\(|AC|=m\)，且
\[
\overrightarrow{BC}=\overrightarrow{AC}-\overrightarrow{AB}=\left(\frac{m}{2},\frac{\sqrt{3}m}{2},0\right),
\]
所以 \(|BC|=m\). 因此 \(\triangle ABC\) 为边长为 \(m\) 的正三角形.

又因为 \(\overrightarrow{AA_{1}}\) 的方向为第三坐标轴方向，而 \(\overrightarrow{AB}\) 和 \(\overrightarrow{AC}\) 的第三个分量均为零，所以 \(AA_{1}\perp AB\)，且 \(AA_{1}\perp AC\). 由此 \(AA_{1}\) 垂直于平面 \(ABC\)，三棱柱为直三棱柱. 其底面是正三角形，所以三棱柱 \(ABC-A_{1}B_{1}C_{1}\) 是正三棱柱.

\item 平面 \(A_{1}ABB_{1}\) 由方向向量 \(\overrightarrow{AB}\) 和 \(\overrightarrow{AA_{1}}\) 张成. 取该平面的一个法向量
\[
\bs{n}=\overrightarrow{AB}\times\overrightarrow{AA_{1}}=\left(-\frac{\sqrt{3}mn}{2},-\frac{mn}{2},0\right),
\]
其长度为 \(|\bs{n}|=mn\). 直线 \(CA_{1}\) 的方向向量可取
\[
\bs{d}=\overrightarrow{CA_{1}}=(-m,0,n),
\]
故 \(|\bs{d}|=\sqrt{m^{2}+n^{2}}\)，且
\[
|\bs{d}\cdot\bs{n}|=\frac{\sqrt{3}m^{2}n}{2}.
\]

设直线 \(CA_{1}\) 与平面 \(A_{1}ABB_{1}\) 所成的角为 \(\theta\)，则
\[
\sin\theta=\frac{|\bs{d}\cdot\bs{n}|}{|\bs{d}|\,|\bs{n}|}=\frac{\sqrt{3}m}{2\sqrt{m^{2}+n^{2}}}.
\]
当 \(m=\sqrt{2}n\) 时，\(\sin\theta=\dfrac{\sqrt{2}}{2}\). 因为 \(0\leqslant\theta\leqslant\dfrac{\pi}{2}\)，所以 \(\theta=\dfrac{\pi}{4}\).
\end{enumerate}
\end{solution}

\begin{problem}
已知函数 \( f(x) = \frac{x^{\frac{1}{3}} - x^{-\frac{1}{3}}}{5} \)，\( g(x) = \frac{x^{\frac{1}{3}} + x^{-\frac{1}{3}}}{5} \).

\begin{enumerate}
\item 证明 \( f(x) \) 是奇函数，并求 \( f(x) \) 的单调区间；
\item 分别计算 \( f(4) - 5f(2)g(2) \) 和 \( f(9) - 5f(3)g(3) \) 的值，由此概括出涉及函数 \( f(x) \) 和 \( g(x) \) 的对所有不等于零的实数 \( x \) 都成立的一个等式，并加以证明.
\end{enumerate}
\end{problem}

\begin{solution}
\begin{enumerate}
\item 函数定义域为非零实数. 对任意 \(x\neq0\)，
\[
f(-x)=\frac{(-x)^{1/3}-(-x)^{-1/3}}{5}=\frac{-x^{1/3}+x^{-1/3}}{5}=-f(x),
\]
所以 \(f(x)\) 是奇函数.

对 \(x\neq0\) 求导，得
\[
f'(x)=\frac{1}{15}x^{-2/3}+\frac{1}{15}x^{-4/3}=\frac{x^{2/3}+1}{15x^{4/3}}>0.
\]
因此 \(f(x)\) 在 \((-\infty,0)\) 和 \((0,+\infty)\) 上都单调递增.

\item 令 \(a=2^{1/3}\)，则
\[
f(4)=\frac{a^{2}-a^{-2}}{5},\qquad 5f(2)g(2)=\frac{(a-a^{-1})(a+a^{-1})}{5}=\frac{a^{2}-a^{-2}}{5},
\]
所以 \(f(4)-5f(2)g(2)=0\).

令 \(b=3^{1/3}\)，则
\[
f(9)=\frac{b^{2}-b^{-2}}{5},\qquad 5f(3)g(3)=\frac{(b-b^{-1})(b+b^{-1})}{5}=\frac{b^{2}-b^{-2}}{5},
\]
所以 \(f(9)-5f(3)g(3)=0\).

由上述计算可猜想，对任意 \(x\neq0\)，有 \(f(x^{2})=5f(x)g(x)\). 直接验证得
\[
f(x^{2})=\frac{x^{2/3}-x^{-2/3}}{5},
\]
而
\[
5f(x)g(x)=5\cdot\frac{(x^{1/3}-x^{-1/3})(x^{1/3}+x^{-1/3})}{25}=\frac{x^{2/3}-x^{-2/3}}{5}.
\]
所以对所有不等于零的实数 \(x\)，恒有 \(f(x^{2})=5f(x)g(x)\).
\end{enumerate}
\end{solution}

\begin{problem}
设 \(F_{1}\)，\(F_{2}\) 分别为椭圆 \(C: \frac{x^{2}}{a^{2}} + \frac{y^{2}}{b^{2}} = 1\)（\(a > 0\)，\(b > 0\)）的左、右两个焦点.

\begin{enumerate}
\item 若椭圆 \(C\) 上的点 \(A\left(1,\frac{3}{2}\right)\) 到 \(F_{1}\)，\(F_{2}\) 两点的距离之和等于 \(4\)，写出椭圆 \(C\) 的方程；
\item 设 \(K\) 是（1）中所得椭圆上的动点，求线段 \(F_{1}K\) 的中点的轨迹方程；
\item 已知椭圆具有性质：若 \(M\)，\(N\) 是椭圆 \(C\) 上关于原点对称的两个点，点 \(P\) 是椭圆上任意一点，当直线 \(PM\)，\(PN\) 的斜率都存在，并记为 \(k_{PM}\)，\(k_{PN}\) 时，那么 \(k_{PM}\) 与 \(k_{PN}\) 之积是与点 \(P\) 位置无关的定值. 试对双曲线 \(\frac{x^{2}}{a^{2}} - \frac{y^{2}}{b^{2}} = 1\) 写出具有类似特性的性质，并加以证明.
\end{enumerate}
\end{problem}

\begin{solution}
\begin{enumerate}
\item 椭圆上任一点到两焦点的距离之和为 \(2a\)，故 \(2a=4\)，得 \(a=2\). 因为点 \(A\left(1,\dfrac{3}{2}\right)\) 在椭圆上，所以
\[
\frac{1}{4}+\frac{(3/2)^{2}}{b^{2}}=1.
\]
解得 \(b^{2}=3\). 因此椭圆 \(C\) 的方程为 \(\dfrac{x^{2}}{4}+\dfrac{y^{2}}{3}=1\).

\item 由 \(a^{2}-b^{2}=1\)，得焦距参数为 \(c=1\)，所以 \(F_{1}=(-1,0)\). 设 \(K(x,y)\)，线段 \(F_{1}K\) 的中点为 \(M(u,v)\)，则
\[
u=\frac{x-1}{2},\qquad v=\frac{y}{2},
\]
即 \(x=2u+1\)，\(y=2v\). 代入椭圆方程，得
\[
\frac{(2u+1)^{2}}{4}+\frac{(2v)^{2}}{3}=1.
\]
所以中点的轨迹方程为
\[
\left(u+\frac{1}{2}\right)^{2}+\frac{4v^{2}}{3}=1.
\]

\item 对双曲线，设 \(M(u,v)\)，\(N(-u,-v)\) 是关于原点对称的两个点，\(P(x,y)\) 是双曲线上任意一点，且两条直线的斜率存在. 则
\[
k_{PM}k_{PN}=\frac{y-v}{x-u}\cdot\frac{y+v}{x+u}=\frac{y^{2}-v^{2}}{x^{2}-u^{2}}.
\]
由 \(P\)，\(M\) 都在双曲线 \(\dfrac{x^{2}}{a^{2}}-\dfrac{y^{2}}{b^{2}}=1\) 上，分别有 \(y^{2}=\dfrac{b^{2}}{a^{2}}x^{2}-b^{2}\) 和 \(v^{2}=\dfrac{b^{2}}{a^{2}}u^{2}-b^{2}\). 相减得
\[
y^{2}-v^{2}=\frac{b^{2}}{a^{2}}(x^{2}-u^{2}).
\]
因此
\[
k_{PM}k_{PN}=\frac{b^{2}}{a^{2}},
\]
这是与点 \(P\) 的位置无关的定值.
\end{enumerate}
\end{solution}

\begin{problem}
在一次人才招聘会上，有 A、B 两家公司分别开出了它们的工资标准：A 公司允诺第一个月工资为 \(1500\text{ 元}\)，以后每年月工资比上一年月工资增加 \(230\text{ 元}\)；B 公司允诺第一年月工资数为 \(2000\text{ 元}\)，以后每年月工资在上一年的月工资基础上递增 \(5\%\)，设某人年初被 A、B 两家公司同时录取. 试问：

\begin{enumerate}
\item 若该人分别在 \(A\) 公司或 \(B\) 公司连续工作 \(n\text{ 年}\)，则他在第 \(n\text{ 年}\)的月工资收入分别是多少？
\item 该人打算连续在一家公司工作 \(10\text{ 年}\)，仅从工资收入总量较多作为应聘的标准（不记其它因素），该人应该选择哪家公司，为什么？
\item 在 \(A\) 公司工作比在 \(B\) 公司工作的月工资收入最多可以多多少元？（精确到 \(1\) 元），并说明理由.
\end{enumerate}
\end{problem}

\begin{solution}
\begin{enumerate}
\item A 公司第 \(n\) 年的月工资构成首项为 \(1500\)、公差为 \(230\) 的等差数列，因此为 \(1500+230(n-1)\) 元. B 公司第 \(n\) 年的月工资构成首项为 \(2000\)、公比为 \(1.05\) 的等比数列，因此为 \(2000(1.05)^{n-1}\) 元.

\item 连续工作 \(10\) 年时，A 公司的月工资总和为
\[
\sum_{k=0}^{9}(1500+230k)=10\cdot1500+230\cdot\frac{9\cdot10}{2}=25350.
\]
B 公司的月工资总和为
\[
\sum_{k=0}^{9}2000(1.05)^k=2000\cdot\frac{1.05^{10}-1}{1.05-1}\approx25155.79.
\]
每年的月工资总和还要同乘 \(12\)，不影响两家公司总量的比较. 因为 \(25350>25155.79\)，所以应选择 A 公司.

\item 设第 \(n\) 年 A 公司与 B 公司月工资之差为
\[
d_n=1500+230(n-1)-2000(1.05)^{n-1}.
\]
则
\[
d_{n+1}-d_n=230-100(1.05)^{n-1}.
\]
又因为 \(1.05^{17}\approx2.2920<2.3<2.4066\approx1.05^{18}\)，所以当 \(n\leqslant18\) 时，\(d_{n+1}-d_n>0\)，当 \(n\geqslant19\) 时，\(d_{n+1}-d_n<0\). 因此 \(d_n\) 在 \(n=19\) 时取得最大值.

此最大值为
\[
d_{19}=1500+230\cdot18-2000(1.05)^{18}\approx826.76\text{ 元}.
\]
精确到 \(1\) 元为 \(827\) 元.
\end{enumerate}
\end{solution}
